\section{Preliminaries}
\label{sec:prelim}

\subsection{Inscribed polytopes}

\begin{cited}{Notation}{{\cite[Section~1]{horvath-langi-2016}}}
  Throughout the paper, by a polytope we mean a convex polytope. The vertex set
  of a polytope $P$ is denoted by $V(P)$. We denote the family of
  $d$-dimensional polytopes, with $n$ vertices and inscribed in the unit sphere
  $\mathbb{S}^{d-1}$, by $\mathcal{P}_{d}(n)$. We denote $d$-dimensional
  volume by $\vol_{d}$, and set
  $v_{d}(n) = \max\{\vol_{d}(P) : P \in \mathcal{P}_{d}(n)\}$.
\end{cited}

\begin{remark}[{\cite[Remark~2.1]{medgyes-maxvol-d4}}]
  \label{rem:attained}
  For $n \geq d + 1$, $v_{d}(n)$ is attained by some
  $P \in \mathcal{P}_{d}(n)$.
\end{remark}

\begin{definition}[{\cite[Definition~2.1, Section~2.2 and Proposition~2.16]{ziegler-1995}}]
  A \emph{facet} of a $d$-polytope is a face of dimension $d-1$, and a
  \emph{ridge} is a face of dimension $d-2$. A polytope is \emph{simplicial} if
  every facet is a simplex. Two polytopes are \emph{combinatorially
  equivalent} if their face lattices are isomorphic; an equivalence class is a
  \emph{combinatorial type}.
\end{definition}

Equivalently, two polytopes are combinatorially equivalent if their
facet--vertex incidence matrices differ only by permutations of rows and
columns \cite[Section~15.1.1]{henk-richter-gebert-ziegler-2017}. For
simplicial polytopes, which is the form in which we use the definition, this
means that a bijection between the vertex sets carries the vertex sets of the
facets of one onto those of the other.

\begin{definition}[{\cite[Definition~1]{horvath-langi-2016}}]
  \label{def:propertyz}
  Let $P \in \mathcal{P}_{d}(n)$ be a $d$-polytope with
  $V(P) = \{p_{1}, p_{2}, \ldots, p_{n}\}$. If for each $i$, there is an open
  set $U_{i} \subset \mathbb{S}^{d-1}$ such that $p_{i} \in U_{i}$, and for any
  $q \in U_{i}$, we have
  \begin{equation*}
    \vol_{d}\bigl(\conv((V(P) \setminus \{p_{i}\}) \cup \{q\})\bigr)
      \leq \vol_{d}(P),
  \end{equation*}
  then we say that $P$ satisfies Property Z.
\end{definition}

\begin{remark}[{\cite{berman-hanes-1970}, \cite[Remark~2.4]{medgyes-maxvol-d4}}]
  \label{rem:propertyz}
  Every $P \in \mathcal{P}_{d}(n)$ with $\vol_{d}(P) = v_{d}(n)$ satisfies
  Property Z.
\end{remark}

\begin{lemma}[{\cite[Lemma~1]{horvath-langi-2016}}]
  \label{lem:simplicial}
  Every $P \in \mathcal{P}_{d}(n)$ satisfying Property Z is simplicial.
\end{lemma}

\subsection{Signed volume and certificates}

For $v_{1}, \ldots, v_{d} \in \mathbb{R}^{d}$, $\det(v_{1}, \ldots, v_{d})$
denotes the determinant of the matrix with columns $v_{1}, \ldots, v_{d}$.

\begin{definition}[{\cite[Definition~3.1]{medgyes-maxvol-d4}}]
  \label{def:signed-volume}
  An \emph{orientation} of a finite set is an ordering of its elements up to
  even permutations. Let $\mathcal{C} = (n, \mathcal{F})$, where
  $\mathcal{F}$ is a finite set of ordered $d$-tuples
  $F = (F_{1}, \ldots, F_{d})$ of distinct elements of $\{1, \ldots, n\}$,
  called facets. Deleting $F_{a}$ gives a ridge, with the orientation of
  $(F_{1}, \ldots, \widehat{F_{a}}, \ldots, F_{d})$ if $a$ is even and the
  opposite one if $a$ is odd. The orientation of $\mathcal{C}$ is
  \emph{coherent} if every ridge lies in exactly two facets and these induce
  opposite orientations on it. A point
  $p = (p_{1}, \ldots, p_{n}) \in (\mathbb{R}^{d})^{n}$ \emph{realises}
  $\mathcal{C}$ if $P = \conv\{p_{1}, \ldots, p_{n}\}$ has the $p_{i}$ as its
  vertices and its facets are exactly the simplices
  $\conv\{p_{j} : j \in F\}$, $F \in \mathcal{F}$. The \emph{signed volume}
  of $p$ is
  \begin{equation}
    \label{eq:signed-volume}
    D(p) = \sum_{F \in \mathcal{F}} \det(p_{F_{1}}, \ldots, p_{F_{d}}).
  \end{equation}
\end{definition}

We write $j \in F$ if $j$ is one of the entries of $F$, and $p_{F}$ for the
set $\{p_{F_{1}}, \ldots, p_{F_{d}}\}$.

\begin{lemma}[{\cite[Lemma~3.2]{medgyes-maxvol-d4}}]
  \label{lem:volume}
  If the orientation of $\mathcal{C}$ is coherent and $p$ realises
  $\mathcal{C}$, then
  \[
    d!\,\vol_{d}(\conv\{p_{1}, \ldots, p_{n}\}) = |D(p)|.
  \]
\end{lemma}

\begin{definition}[{\cite[Definition~3.4]{medgyes-maxvol-d4}}]
  \label{def:certificate}
  Let $x_{i} = (x_{i,1}, \ldots, x_{i,d})$ be the coordinates of $p_{i}$, and
  $v(x)$ the vector of all monomials of degree at most $2$ in the $x_{i,a}$. A
  \emph{certificate} for $\mathcal{C}$ and $\beta \in \mathbb{Q}$ is a
  symmetric positive definite matrix $Q$ with rational entries together with
  polynomials $h_{1}, \ldots, h_{n}$ with rational coefficients such that
  \begin{equation}
    \label{eq:certificate}
    \beta - D(x) = v(x)^{\mathsf{T}} Q\, v(x)
      + \sum_{i=1}^{n} h_{i}(x)\bigl(|x_{i}|^{2} - 1\bigr)
  \end{equation}
  holds identically as polynomials in the $x_{i,a}$.
\end{definition}

\begin{proposition}[{\cite[Proposition~3.5]{medgyes-maxvol-d4}}]
  \label{prop:certificate}
  Suppose the orientation of $\mathcal{C}$ is coherent and a certificate for
  $\mathcal{C}$ and $\beta$ exists. Then $|D(p)| \leq \beta$ for every
  $p \in (\mathbb{S}^{d-1})^{n}$, and every polytope inscribed in
  $\mathbb{S}^{d-1}$ that realises $\mathcal{C}$ has volume at most
  $\beta/d!$.
\end{proposition}

The first conclusion of Proposition~\ref{prop:certificate} concerns every
configuration on the sphere, convex or not. A certificate therefore cannot
exist for $\beta$ below the maximum of $|D|$ on
$(\mathbb{S}^{d-1})^{n}$, even if that maximum is attained only by
configurations that realise no polytope of the type.
Section~\ref{sec:convexity} removes this restriction.

\subsection{The join of two pentagons}

We use the following results of \cite{medgyes-maxvol-d4-n9} with
$m = l = 5$. For $m, l \geq 3$, $\mathcal{J}_{m,l}$ denotes the join of an
$m$-cycle and an $l$-cycle \cite[Definition~3.1]{medgyes-maxvol-d4-n9}, a
coherently oriented type with $m + l$ vertices
\cite[Lemma~3.3]{medgyes-maxvol-d4-n9}.

\begin{lemma}[{\cite[Lemma~3.2]{medgyes-maxvol-d4-n9}}]
  \label{lem:freesum}
  Let $m, l \geq 3$, let $u_{1}, u_{2}, u_{3}, u_{4}$ be an orthonormal basis of
  $\mathbb{R}^{4}$, and for $i \in \mathbb{Z}_{m}$, $j \in \mathbb{Z}_{l}$ let
  \[
    a_{i} = \cos\tfrac{2\pi i}{m}\,u_{1} + \sin\tfrac{2\pi i}{m}\,u_{2}, \qquad
    b_{j} = \cos\tfrac{2\pi j}{l}\,u_{3} + \sin\tfrac{2\pi j}{l}\,u_{4}.
  \]
  Then $P = \conv\{a_{i}, b_{j}\} \in \mathcal{P}_{4}(m + l)$, and the point
  $p = (a_{1}, \ldots, a_{m}, b_{1}, \ldots, b_{l})$ realises
  $\mathcal{J}_{m,l}$.
\end{lemma}

\begin{corollary}[{\cite[Corollary~5.2]{medgyes-maxvol-d4-n9}}]
  \label{cor:pair}
  Let $m, l \in \{3, 4, 5\}$. If $P \in \mathcal{P}_{4}(m+l)$ is
  combinatorially equivalent to the polytope of Lemma~\ref{lem:freesum}, then
  \[
    \vol_{4}(P) \leq \frac{1}{24}\, m\sin\frac{2\pi}{m}\; l\sin\frac{2\pi}{l},
  \]
  with equality if and only if $P$ is the polytope of Lemma~\ref{lem:freesum}
  for some orthonormal basis $u_{1}, \ldots, u_{4}$, that is, up to an
  orthogonal transformation, the convex hull of a regular $m$-gon and a
  regular $l$-gon inscribed in $\mathbb{S}^{3}$ in orthogonal $2$-planes
  through the origin.
\end{corollary}

For $m = l = 5$ the bound is $V^{*} = \frac{25}{24}\sin^{2}\frac{2\pi}{5}$.
Since $\sin^{2}\frac{2\pi}{5} = \frac{1}{16}(10 + 2\sqrt{5})$,
\begin{equation}
  \label{eq:vstar}
  24\,V^{*} = \frac{25}{16}\bigl(10 + 2\sqrt{5}\bigr), \qquad
  V^{*} = \frac{25}{384}\bigl(10 + 2\sqrt{5}\bigr) = 0.942196\ldots.
\end{equation}
